\begin{figure}[H]
\centering
\begin{tikzpicture}[ancho/.style={minimum width=3.3cm}]
    \node[documento, ancho] (current) at (0,0) {Contenido\\actual};
    \node[servicio, ancho] (snapshot) at (5.3,0) {Servicio de\\respaldos};
    \node[documento, ancho] (copy) at (10.6,0) {Carpeta de\\respaldo};
    \node[datos, minimum width=3cm] (metadb) at (5.3,-3) {SQLite\\metadatos};
    \node[documento, ancho] (archive) at (10.6,-3) {Archivo\\transferible};

    \draw[flecha] (current) -- node[etq, above=3pt] {copia} (snapshot);
    \draw[flecha] (snapshot) -- node[etq, above=3pt] {guarda\\archivos} (copy);
    \draw[flecha] (snapshot) -- node[etq, right=3pt] {registra\\autor y fecha} (metadb);
    \draw[flecha, {Stealth[length=2.8mm, width=2mm]}-{Stealth[length=2.8mm, width=2mm]}] (copy) -- node[etq, right=3pt] {exporta\\o importa} (archive);
    \draw[flecha] (copy.north) -- ++(0,0.75) -| node[etq, pos=0.25, above=2pt] {restaura el contenido} (current.north);
\end{tikzpicture}
\caption[Ciclo de respaldos del CMS]{Ciclo de creación, transferencia y restauración de respaldos del contenido}
\label{fig:cms_ciclo_respaldos}
\end{figure}
